\section{Separating execution failure from scientific evidence}

Let $R_q$ be all structurally valid receipts for request $q$. Partition them into
completed outcomes $R_q^+$, valid scientific negatives $R_q^-$, uncheckable
attempts $R_q^?$ and execution failures $R_q^f$. Every attempt should be retained
in an execution ledger, because omission can hide reliability failures. Only the
first two sets are candidates for a scientific evidence ledger, and both remain
subject to the frozen adjudication $F$.

\begin{proposition}[Failure non-contamination]
If $Z(r)$ accepts only terminals in $R_q^+\cup R_q^-$, then no receipt in
$R_q^f$ can enter the admitted scientific-evidence set through
$\operatorname{Admit}$.
\end{proposition}

\begin{proof}
For $r\in R_q^f$, the closed terminal type assigns $z=\failure$. By definition,
$Z(r)$ is false for this terminal. Therefore the conjunction defining
$\operatorname{Admit}$ is false independently of all other predicates.
\end{proof}

This is a contract theorem, not empirical evidence that an implementation always
enforces the contract. A hostile implementation can misclassify a network failure
as a completed empty result, bypass the typed constructor or mutate stored bytes.
The adversarial protocol in Section~\ref{sec:attacks} must therefore test the
implementation boundary.

\begin{proposition}[No silent row deletion]
If every accepted request identifier must terminate in exactly one sealed receipt
or one explicit missing-receipt alarm, then a campaign summary cannot convert an
execution failure into an absent observation without violating the coverage
predicate.
\end{proposition}

\begin{proof}
The coverage predicate compares the frozen request set with the disjoint union of
sealed receipts and explicit alarms. Removing a failed attempt from both sets makes
the union smaller than the request set, so coverage fails.
\end{proof}

Valid negative results require the opposite treatment: they must remain in the
scientific ledger. Preregistration and Registered Reports address outcome-dependent
analysis and publication incentives \citep{nosek2018preregistration,chambers2015registered}.
The execution contract supports that separation by binding the protocol and
terminal schema before outcome access, but cannot determine whether the registered
scientific design is meaningful.

Repair verification is also distinct from scientific replication. Rerunning a
fixed case after correcting a digest, parser or environment defect can establish
that the repaired implementation meets the same contract. Unless the scientific
protocol prospectively defines the repair run as a new independent unit, it cannot
be counted as additional corroboration.

